% ==========================================
% BAB VI EVALUASI
% ==========================================
\cleardoublepage
\chapter{EVALUASI}
\label{chap:evaluasi}
Bab ini memaparkan proses evaluasi terhadap sistem \textit{Warehouse Management Controller} (WMC) dan Subsistem Lokalisasi yang telah diimplementasikan. Evaluasi bertujuan untuk mengukur sejauh mana sistem yang dikembangkan memenuhi spesifikasi yang telah ditetapkan pada dokumen B200. Untuk WMC, cakupan evaluasi meliputi akurasi pemenuhan tugas dan ketahanan sistem di bawah beban (\textit{stress test}), sedangkan untuk Subsistem Lokalisasi, cakupan evaluasi meliputi akurasi pemetaan posisi dan deteksi orientasi. Bab ini disusun menjadi tiga sub-bab: Subbab \ref{sec:metode-evaluasi} memaparkan metode evaluasi yang digunakan, Subbab \ref{sec:hasil-verifikasi} menyajikan hasil verifikasi dari setiap metode tersebut, dan Subbab \ref{sec:pembahasan-evaluasi} membahas hasil evaluasi secara keseluruhan sebagai dasar penilaian kelayakan sistem.

\section{Metode Evaluasi}
\label{sec:metode-evaluasi}

\subsection{Metode Evaluasi Sistem WMC}
Evaluasi sistem WMC dirancang mengacu pada rencana verifikasi yang tertuang dalam dokumen B200, dengan memfokuskan penilaian pada komponen-komponen yang secara langsung dipengaruhi oleh WMC sebagai lapisan perangkat lunak pengendali gudang. Mengingat peran WMC sebagai sistem \textit{backend} yang mengatur manajemen inventori, koordinasi tugas robot, serta antarmuka pengguna, maka terdapat dua kriteria evaluasi utama yang akan diukur, yaitu \textit{Task Fulfilment Accuracy} dan \textit{Stress Test}. Setiap kriteria dievaluasi melalui metode pengujian yang berbeda namun saling melengkapi untuk memperoleh gambaran yang komprehensif mengenai kualitas sistem secara keseluruhan.

\paragraph{a. Pengujian Task Fulfilment Accuracy --- Verifikasi Konsistensi Data}
Pengujian akurasi pemenuhan tugas bertujuan untuk memverifikasi bahwa sistem WMC mampu menghasilkan dan mengelola perintah tugas dengan tingkat akurasi sebesar 100\%, tanpa terjadi kesalahan lokasi rak (\textit{misplacement}) dan tanpa ketidaksesuaian antara kondisi fisik dengan data inventori yang tersimpan. Dari sisi WMC, kriteria ini mencerminkan kualitas logika \textit{Fleet Order Service} dalam menentukan lokasi target, serta ketepatan \textit{Goods Detail Service} dalam memperbarui status dan posisi barang setelah setiap tugas diselesaikan.

Metode pengujian dilakukan dengan membandingkan data yang tercatat dalam basis data WMC terhadap kondisi fisik aktual setelah serangkaian tugas pengambilan dan penempatan barang dieksekusi. Pengujian dilaksanakan dengan menjalankan 8 \textit{order storing} dan 8 \textit{order picking} (total 16 \textit{task}) menuju 8 rak \textit{storage} yang berbeda, dengan robot memulai pergerakan dari 4 sel \textit{idle} yang berbeda secara bergantian. Setiap \textit{order} yang dibuat melalui antarmuka WMC dicatat beserta lokasi tujuan yang ditetapkan oleh sistem. Setelah robot menyelesaikan tugas, data lokasi barang yang diperbarui oleh sistem dibandingkan dengan lokasi fisik aktual di gudang. Tingkat akurasi dihitung menggunakan formula berikut:
\begin{equation}
\text{Task Fulfilment Accuracy} = \frac{\text{Jumlah Tugas Berhasil Tanpa Kesalahan}}{\text{Jumlah Total Tugas}} \times 100\%
\label{eq:tfa}
\end{equation}
Sistem dinyatakan lulus pengujian apabila seluruh tugas yang dieksekusi memberikan hasil yang sesuai antara perintah sistem dengan kondisi fisik, sehingga akurasi mencapai 100\%.

\paragraph{b. Pengujian Stress Test --- Pengukuran Ketahanan Sistem di Bawah Beban}
Pengujian \textit{stress test} dilakukan untuk mengukur ketahanan dan performa sistem WMC ketika menerima beban permintaan yang tinggi secara bersamaan pada setiap \textit{endpoint} utama. Sistem dinyatakan memenuhi kriteria apabila seluruh \textit{endpoint} yang diuji tidak mengalami kegagalan (\textit{error rate} 0\%) dan tetap responsif di bawah beban, yang diukur melalui waktu respons rata-rata, persentil ke-95 (p95), dan persentil ke-99 (p99).

Metode pengujian dilakukan dengan mengirimkan sejumlah \textit{request} secara simultan ke setiap \textit{endpoint} pada \textit{Authentication Service} dan \textit{Goods Service} menggunakan alat pengujian beban (\textit{load testing tool}). Setiap \textit{endpoint} diuji secara terpisah dengan mencatat waktu respons rata-rata, p95, p99, \textit{error rate}, serta \textit{throughput} (jumlah \textit{request} yang berhasil diproses per detik). Hasil pengujian kemudian dibandingkan terhadap kriteria keberhasilan, yaitu tidak ada kegagalan \textit{request} selama pengujian berlangsung.

Rangkuman seluruh kriteria evaluasi beserta parameter target dan metode pengujiannya disajikan pada Tabel \ref{tab:rangkuman-evaluasi}.

\begin{table}[h]
  \centering
  \caption{Rangkuman kriteria dan metode evaluasi WMC}
  \label{tab:rangkuman-evaluasi}
  \begin{tabularx}{\textwidth}{ | c | l | X | X | }
    \hline
    \textbf{No} & \textbf{Kriteria Evaluasi} & \textbf{Parameter Target (B200)} & \textbf{Metode Pengujian} \\
    \hline
    1 & Task Fulfilment Accuracy & Akurasi 100\%; tanpa \textit{misplacement}; konsistensi data fisik dan inventori & Perbandingan data basis data WMC terhadap kondisi fisik aktual pasca eksekusi tugas \\
    \hline
    2 & Stress Test & \textit{Error rate} 0\% pada seluruh \textit{endpoint} di bawah beban & Pengujian beban simultan pada \textit{endpoint} utama dengan alat pengujian beban \\
    \hline
  \end{tabularx}
\end{table}

Kedua metode evaluasi tersebut dirancang agar dapat dilaksanakan secara bertahap dan sistematis. Pengujian \textit{stress test} dapat dilakukan secara mandiri dalam lingkungan pengembangan, sementara pengujian akurasi pemenuhan tugas memerlukan integrasi penuh antara sistem WMC dengan subsistem robot dan \textit{Fleet Controller} dalam skenario operasional terintegrasi.

\subsection{Metode Evaluasi Sistem Lokalisasi}
Verifikasi Subsistem Lokalisasi dilakukan dengan menempatkan \textit{marker} Tag36h11 AprilTag ($8\times8$ cm) pada seluruh sel \textit{grid} lantai gudang berukuran $8\times4$ (32 sel), dengan posisi dan orientasi yang telah diketahui secara pasti sebagai \textit{ground truth}. Pengujian orientasi dilakukan pada empat kelipatan sudut, yaitu $0^\circ$, $90^\circ$, $180^\circ$, dan $270^\circ$, sesuai dengan resolusi deteksi \textit{heading} yang digunakan oleh sistem.

Pengujian dilaksanakan dengan menempatkan banyak \textit{marker} AprilTag secara bersamaan pada seluruh 32 sel \textit{grid} untuk memastikan bahwa pembacaan posisi berhasil dilakukan secara konsisten pada setiap sel, tidak hanya pada sebagian titik sampel. \textit{Payload} MQTT yang dipublikasikan oleh subsistem kemudian dibandingkan terhadap \textit{ground truth} untuk dua properti: posisi sel \textit{grid} (kolom, baris) dan sudut orientasi (\textit{heading}). Target keberhasilan adalah akurasi 100\% pada kedua properti tersebut di seluruh sel \textit{grid} dan seluruh sudut orientasi yang diuji. Tabel pengujian yang disajikan pada Subbab \ref{subsec:hasil-lokalisasi} menampilkan sebagian titik uji sebagai representasi dari keseluruhan pengujian yang telah dilakukan pada seluruh 32 sel \textit{grid}.

\section{Hasil Verifikasi}
\label{sec:hasil-verifikasi}
Sub-bab ini menyajikan hasil pelaksanaan setiap metode evaluasi yang telah direncanakan pada Subbab \ref{sec:metode-evaluasi}, dipisahkan antara hasil verifikasi Sistem WMC pada Subbab \ref{subsec:hasil-wmc} dan hasil verifikasi Sistem Lokalisasi pada Subbab \ref{subsec:hasil-lokalisasi}.

\subsection{Hasil Verifikasi Sistem WMC}
\label{subsec:hasil-wmc}

\subsubsection*{VI.2.1.1 Hasil Pengujian Task Fulfilment Accuracy}
Pengujian akurasi pemenuhan tugas dilaksanakan dengan mengeksekusi 8 tugas \textit{storing} dan 8 tugas \textit{picking} melalui sistem WMC menuju 8 rak \textit{storage} yang berbeda, dengan robot memulai pergerakan dari 4 sel \textit{idle} yang berbeda secara bergantian, kemudian membandingkan data yang tercatat dalam basis data dengan kondisi fisik aktual di gudang. Setiap \textit{task} diperiksa untuk memastikan tidak terjadi kesalahan lokasi rak (\textit{misplacement}) dan tidak ada ketidaksesuaian antara kondisi fisik dengan data inventori. Hasil pengujian disajikan pada Tabel \ref{tab:tfa}.

\begin{table}[h]
  \centering
  \caption{Hasil pengujian Task Fulfilment Accuracy}
  \label{tab:tfa}
  \begin{tabular}{ | l | c | c | c | }
    \hline
    \textbf{Tipe Task} & \textbf{Jumlah Task Dieksekusi} & \textbf{Task Tanpa Kesalahan Lokasi} & \textbf{Akurasi (\%)} \\
    \hline
    Storing (penempatan barang) & 8 & 8 & 100\% \\
    \hline
    Picking (pengambilan barang) & 8 & 8 & 100\% \\
    \hline
    \textbf{Total keseluruhan} & \textbf{16} & \textbf{16} & \textbf{100\%} \\
    \hline
  \end{tabular}
\end{table}

Seluruh 16 \textit{task} yang dieksekusi menuju 8 rak \textit{storage} dengan titik awal dari 4 sel \textit{idle} yang berbeda menunjukkan kesesuaian penuh antara lokasi target yang ditetapkan oleh sistem WMC dengan kondisi fisik aktual di gudang. Tidak terdapat satu pun kejadian \textit{misplacement} maupun ketidaksesuaian data inventori, sehingga tingkat \textit{Task Fulfilment Accuracy} mencapai 100\%, memenuhi target spesifikasi B200.

\subsubsection*{VI.2.1.2 Hasil Pengujian Stress Test}
Pengujian \textit{stress test} dilaksanakan dengan mengirimkan \textit{request} secara simultan ke sebelas \textit{endpoint} utama pada \textit{Authentication Service} dan \textit{Goods Service}. Setiap \textit{endpoint} diukur waktu respons rata-rata, p95, p99, \textit{error rate}, dan \textit{throughput} selama pengujian berlangsung. Hasil pengujian dirangkum pada Tabel \ref{tab:stress-test}.

\begin{table}[h]
  \centering
  \caption{Hasil pengujian stress test \textit{endpoint} WMC}
  \label{tab:stress-test}
  \begin{tabularx}{\textwidth}{ | X | c | c | c | c | c | c | }
    \hline
    \textbf{Endpoint} & \textbf{Avg} & \textbf{p(95)} & \textbf{p(99)} & \textbf{Error Rate} & \textbf{Throughput} & \textbf{Status} \\
    \hline
    POST /auth/login & 6,12 ms & 13,49 ms & 52,36 ms & 0,00\% & 17,6 req/s & Pass All \\
    \hline
    GET /items (Goods Service) & 5,06 ms & 10,37 ms & 120,6 ms & 0,00\% & 99,4 req/s & Pass All \\
    \hline
    POST /items (Goods Service) & 10,58 ms & 18,2 ms & 46,24 ms & 0,00\% & 19,8 req/s & Pass All \\
    \hline
    GET /orders (Blob Service) & 5,12 ms & 11,13 ms & 21,96 ms & 0,00\% & 19,9 req/s & Pass All \\
    \hline
    GET /apriltag-mappings/status & 5,54 ms & 17,06 ms & 26,28 ms & 0,00\% & 19,9 req/s & Pass All \\
    \hline
    POST /apriltag-mappings & 9 ms & 15,26 ms & 20,87 ms & 0,00\% & 19,8 req/s & Pass All \\
    \hline
    GET /orders & 5,93 ms & 11,74 ms & 28,57 ms & 0,00\% & 19,9 req/s & Pass All \\
    \hline
    POST /orders & 8,71 ms & 15,29 ms & 24,63 ms & 0,00\% & 19,6 req/s & Pass All \\
    \hline
    GET /apriltag-mappings/item/:itemId & 6,21 ms & 13,44 ms & 22,87 ms & 0,00\% & 19,83 req/s & Pass All \\
    \hline
    POST /apriltag-mappings/scan & 8,45 ms & 14,92 ms & 21,33 ms & 0,00\% & 19,79 req/s & Pass All \\
    \hline
    PATCH /apriltag-mappings/:mappingId/status & 10,13 ms & 17,58 ms & 24,61 ms & 0,00\% & 19,71 req/s & Pass All \\
    \hline
  \end{tabularx}
\end{table}

Seluruh sebelas \textit{endpoint} yang diuji menunjukkan \textit{error rate} sebesar 0,00\%, yang berarti tidak ada \textit{request} yang gagal diproses selama pengujian berlangsung, sehingga memenuhi kriteria keberhasilan yang ditetapkan. Waktu respons p99 tertinggi tercatat pada \textit{endpoint} \texttt{GET /items} sebesar 120,6 ms, sementara \textit{endpoint} lainnya secara konsisten berada di bawah 55 ms pada persentil ke-99.

\subsection{Hasil Verifikasi Sistem Lokalisasi}
\label{subsec:hasil-lokalisasi}

\begin{table}[h]
  \centering
  \caption{Contoh hasil verifikasi posisi Subsistem Lokalisasi (representasi dari pengujian pada seluruh 32 sel grid $8\times4$)}
  \label{tab:verifikasi-posisi}
  \begin{tabular}{ | c | c | c | }
    \hline
    \textbf{Tag ID} & \textbf{Sel Grid yang Diharapkan (c, r)} & \textbf{Status} \\
    \hline
    562 & (0, 1) & \checkmark Lulus \\
    \hline
    584 & (0, 2) & \checkmark Lulus \\
    \hline
    577 & (1, 2) & \checkmark Lulus \\
    \hline
    585 & (3, 1) & \checkmark Lulus \\
    \hline
    586 & (5, 1) & \checkmark Lulus \\
    \hline
    578 & (2, 3) & \checkmark Lulus \\
    \hline
  \end{tabular}
\end{table}

\begin{table}[h]
  \centering
  \caption{Contoh hasil verifikasi orientasi Subsistem Lokalisasi (kelipatan $90^\circ$)}
  \label{tab:verifikasi-orientasi}
  \begin{tabular}{ | c | c | c | }
    \hline
    \textbf{Tag ID} & \textbf{Heading yang Diharapkan} & \textbf{Status} \\
    \hline
    584 & $0^\circ$ & \checkmark Lulus \\
    \hline
    562 & $90^\circ$ & \checkmark Lulus \\
    \hline
    578 & $180^\circ$ & \checkmark Lulus \\
    \hline
    577 & $270^\circ$ & \checkmark Lulus \\
    \hline
  \end{tabular}
\end{table}

Pengujian posisi dilaksanakan dengan menempatkan AprilTag pada seluruh 32 sel \textit{grid} $8\times4$ secara bersamaan, dan seluruh sel berhasil terdeteksi dan dipetakan ke koordinat \textit{grid} yang sesuai dengan \textit{ground truth}, mencapai akurasi 100\%. Tabel \ref{tab:verifikasi-posisi} menampilkan enam titik uji sebagai contoh representatif dari keseluruhan pengujian tersebut. Pengujian orientasi dilaksanakan pada empat sudut \textit{heading} ($0^\circ$, $90^\circ$, $180^\circ$, $270^\circ$) dan seluruhnya terdeteksi dengan benar, sebagaimana dicontohkan pada Tabel \ref{tab:verifikasi-orientasi}.

\section{Pembahasan Hasil Evaluasi}
\label{sec:pembahasan-evaluasi}
Berdasarkan hasil verifikasi pada Subbab \ref{sec:hasil-verifikasi}, seluruh kriteria evaluasi yang ditetapkan pada dokumen B200 untuk sistem WMC dan Subsistem Lokalisasi telah terpenuhi. Pada sisi WMC, pengujian \textit{Task Fulfilment Accuracy} menunjukkan akurasi 100\% dari 16 \textit{task} (8 \textit{storing} dan 8 \textit{picking}) yang dieksekusi menuju 8 rak \textit{storage} berbeda dengan titik awal dari 4 sel \textit{idle} yang berbeda, tanpa satu pun kejadian \textit{misplacement} atau ketidaksesuaian data inventori, yang mengonfirmasi bahwa logika \textit{Fleet Order Service} dalam menentukan lokasi target serta mekanisme pembaruan data pada \textit{Goods Detail Service} berjalan dengan benar dan andal. Pengujian \textit{stress test} pada sebelas \textit{endpoint} utama \textit{Authentication Service} dan \textit{Goods Service} menunjukkan \textit{error rate} 0,00\% di seluruh \textit{endpoint}, dengan waktu respons p99 tertinggi sebesar 120,6 ms pada \textit{endpoint} \texttt{GET /items}, sementara \textit{endpoint} lainnya secara konsisten berada di bawah 55 ms pada persentil ke-99. Hasil ini mengindikasikan bahwa sistem WMC tetap responsif dan tidak mengalami kegagalan pemrosesan permintaan sekalipun menerima beban tinggi secara bersamaan.

Pada sisi Subsistem Lokalisasi, verifikasi yang dilakukan pada seluruh 32 sel \textit{grid} lantai gudang berukuran $8\times4$ menunjukkan akurasi 100\% untuk pemetaan posisi, dan verifikasi pada seluruh empat kelipatan sudut orientasi ($0^\circ$, $90^\circ$, $180^\circ$, $270^\circ$) juga menunjukkan akurasi 100\% untuk deteksi \textit{heading}. Hasil ini mengonfirmasi bahwa \textit{pipeline} Subsistem Lokalisasi---mulai dari koreksi distorsi \textit{fisheye}, transformasi \textit{homography}, hingga \textit{snap} ke sel \textit{grid} diskrit---berfungsi dengan benar dan mampu melakukan pemetaan posisi serta deteksi orientasi secara andal di seluruh area lantai gudang.

Perlu dicatat bahwa data hasil \textit{stress test} yang disajikan pada Tabel \ref{tab:stress-test} merupakan hasil pengujian awal dan masih bersifat sementara, sehingga berpotensi diperbarui seiring dengan pengujian lanjutan pada kondisi beban yang lebih tinggi atau skenario operasional yang lebih bervariasi. Demikian pula, tabel verifikasi posisi dan orientasi Subsistem Lokalisasi (Tabel \ref{tab:verifikasi-posisi} dan Tabel \ref{tab:verifikasi-orientasi}) hanya menampilkan sebagian titik uji sebagai representasi dari keseluruhan 32 sel \textit{grid} dan seluruh sudut orientasi yang telah diverifikasi secara penuh.

Secara keseluruhan, hasil evaluasi menunjukkan bahwa sistem \textit{Warehouse Management Controller} (WMC) dan Subsistem Lokalisasi yang dikembangkan telah memenuhi seluruh kriteria evaluasi yang ditetapkan. Ketiga aspek yang diuji, yaitu akurasi pemenuhan tugas, ketahanan sistem di bawah beban (\textit{stress test}), serta akurasi lokalisasi posisi dan orientasi, semuanya menunjukkan hasil yang memenuhi atau melampaui nilai target minimum yang disyaratkan dalam spesifikasi B200. Hal ini mengonfirmasi bahwa WMC dan Subsistem Lokalisasi yang diimplementasikan layak digunakan sebagai komponen inti dalam sistem \textit{Smart Warehouse} Paragon Corp.
